% ============ 第4章 §4.3 Calderón--Zygmund 分解与弱 (1,1) 有界性 (原书 pp.82-88) ============

\section{Calderón--Zygmund 分解与弱 (1,1) 有界性}\label{sec:4.3}
\index{0symbolEjf@$E_{j}f$}

\index{0symbolMdyadic@$M_{\dyadic}f$}

\index{C@Calder\'on--Zygmund 分解}

\index{E@二进制极大算子}

\index{E@二进制区间}

\index{L@Lebesgue 微分定理}














可测函数 $f: \mathbb{R} \to \mathbb{C}$ 的 $L^{2}$ 范数由
\[
\|f\|_{L^{2}} = \left(\int_{\mathbb{R}}|f(x)|^{2}\dd x\right)^{\frac{1}{2}}
\]
给出, 同样地 $L^{1}$ 范数由
\[
\|f\|_{L^{1}} = \int_{\mathbb{R}}|f(x)|\dd x
\]
给出. 前一节证明了 Hilbert 变换 $H$ 的 $L^{2}$ 有界性, 这里想考虑 $H$ 的 $L^{1}$ 有界性. 遗憾的是（练习题~\ref{exer:4.10} 参照）,
\[
H\chi_{[-1,1]}(x) = -\log\left|\frac{1-x}{1+x}\right| = O(x^{-1}) \quad (|x| \to \infty),
\]
所以 $H$ 不是 $L^{1}$ 有界的. Hilbert 变换定义为
\[
Hf(x) = \lim_{\varepsilon \downarrow 0}\int_{|y| > \varepsilon}\frac{f(x-y)}{y}\dd y,
\]
但回顾: 若 $f \in L^{2}(\mathbb{R})$, 则在 $L^{2}$ 的拓扑下收敛, 且 $H$ 是（由 Cotlar 引理）$L^{2}$ 有界的. 关于 Hilbert 变换的显著性质仅此而已, 剩下的只是利用 Hilbert 变换大致上是由与 $x^{-1}$ 的卷积给出这一点来推进本节的讨论.

作为考察 Hilbert 变换性质的起点, 考察对 $L^{1}$ 函数的输入有什么样的性质. 此时, 使用 1950 年代构思的 Calderón--Zygmund 分解.

这里定义在调和分析中起重要作用的二进极大算子与二进区间.

\begin{definition}[二进区间]\label{def:4.1}
对 $j, k \in \mathbb{Z}$, 区间 $I_{jk}$ 由
\begin{equation}\tag{4.9}
I_{jk} \equiv \left[\frac{k}{2^{j}}, \frac{k+1}{2^{j}}\right)
\label{eq:4.9}
\end{equation}
定义. 如此用 $j, k \in \mathbb{Z}$ 表为 $I_{jk}$ 的集合称为二进区间. 全体二进区间组成的集合记为 $\mathcal{D} = \mathcal{D}(\mathbb{R}) \equiv \{I_{jk} : j \in \mathbb{Z}, k \in \mathbb{Z}\}$. $j$ 世代的区间是指这些区间中属于
\begin{equation}\tag{4.10}
\mathcal{D}_{j} = \mathcal{D}_{j}(\mathbb{R}) \equiv \{Q \in \mathcal{D}(\mathbb{R}) : \ell(Q) = 2^{-j}\}
\label{eq:4.10}
\end{equation}
的全体.
\end{definition}

对二进区间 Lebesgue 微分定理也成立. 即, 对每个 $x \in \mathbb{R}$, 用 $Q_{j,x}$ 表示满足 $x \in Q_{j,x} \in \mathcal{D}_{j}(\mathbb{R})$ 的唯一区间, 则对任意局部可积函数 $f$, 几乎处处
\[
f(x) = \lim_{j \to \infty}2^{j}\int_{Q_{j,x}}f(y)\dd y
\]
成立. 实际上, 对连续函数这件事是显然的. 这个关系也称为 Lebesgue 微分定理.

下面的引理几乎显然, 省略证明（定理~\ref{thm:1.7} 也参照）.

\begin{lemma}[二进区间的嵌套结构]\label{lem:4.1}
两个二进区间相交时, 一方包含另一方.
\end{lemma}

定义了二进区间, 下面定义二进极大算子. 对区间 $I$, 用 $m_{I}(f)$ 表示 $f$ 在 $I$ 上的平均.

\begin{definition}[二进极大算子]\label{def:4.2}
对 $\mathbb{R}$ 上的局部可积函数 $f$,
\begin{align*}
E_{j}f(x) &\equiv E_{j}[f](x) \equiv \sum_{Q \in \mathcal{D}_{j}(\mathbb{R})}m_{Q}(f) \cdot \chi_{Q}(x), \\
M_{\dyadic}f(x) &\equiv \sup_{j \in \mathbb{Z}}E_{j}[|f|](x),
\end{align*}
并称 $M_{\dyadic}$ 为二进极大算子. $E_{j}f(x)$ 是局部可积函数 $f$ 在包含 $x$ 的唯一 $j$ 世代区间 $Q_{j,x} \in \mathcal{D}_{j}(\mathbb{R})$ 上的平均.
\end{definition}

$E_{j}$ 是线性的, 但 $M_{\dyadic}$ 不是线性的.

对一切 $j \in \mathbb{Z}$ 与 $x \in \mathbb{R}$,
\[
|E_{j}f(x)| = \sum_{Q \in \mathcal{D}_{j}}|m_{Q}(f)\chi_{Q}(x)| \leq 2^{j}\|f\|_{L^{1}}\sum_{Q \in \mathcal{D}_{j}}\chi_{Q}(x) = 2^{j}\|f\|_{L^{1}}
\]
可知不等式 $|E_{j}f(x)| \leq 2^{j}\|f\|_{L^{1}}$ 成立, 所以下面的引理是显然的.

\begin{lemma}\label{lem:4.2}
对任意 $f \in L^{1}(\mathbb{R})$ 与 $x \in \mathbb{R}$, $\lim_{j \to -\infty}E_{j}f(x) = 0$ 成立.
\end{lemma}

下面的定理是 Hilbert 变换的弱 (1,1) 有界性的证明的关键的函数分解, 是调和分析学的重要工具之一. 用 $\coprod$ 表示不相交的并.

\begin{theorem}[Calderón--Zygmund 区间的构造]\label{thm:4.2}
对可积函数 $f$ 与 $\lambda > 0$, 存在互不相交的二进区间的集合 $\{I_{j}\}_{j \in J} \subset \mathcal{D}$, 使得
\begin{equation}\tag{4.11}
\{x \in \mathbb{R} : M_{\dyadic}f(x) > \lambda\} = \coprod_{j \in J}I_{j}, \quad \lambda < m_{I_{j}}(|f|) \leq 2\lambda
\label{eq:4.11}
\end{equation}
成立.
\end{theorem}

\begin{proof}
为简化记号, 令 $E_{\lambda} \equiv \{x \in \mathbb{R} : M_{\dyadic}f(x) > \lambda\}$. 由二进极大算子的定义式
\[
M_{\dyadic}f(x) = \sup_{Q \in \mathcal{D}}m_{Q}(|f|)\chi_{Q}(x),
\]
特别地由这里出现的 sup 的性质, $E_{\lambda} = \bigcup_{j \in \mathbb{Z}}\{E_{j}[|f|] > \lambda\}$ 成立. 由引理~\ref{lem:4.2}, 对每个 $x \in E_{\lambda}$, 使 $E_{j}[|f|](x) > \lambda$, 即
\[
x \in \{y \in \mathbb{R} : E_{j}[|f|](y) > \lambda\}
\]
成立的 $j$ 存在下界. 即, 存在 $j^{*}$, 使得对一切 $j < j^{*}$, $E_{j}[|f|](x) \leq \lambda$ 成立. 因此, $E_{\lambda}$ 可以分割为
\[
E_{\lambda} = \coprod_{j \in \mathbb{Z}}\{x \in \mathbb{R} : E_{j}[|f|](x) > \lambda, \text{ 且 } k < j \text{ 时 } E_{k}[|f|](x) \leq \lambda\}.
\]
这里, 设 $j_{M} \in \mathbb{N} \cup \{\infty\}$,
\begin{equation}\tag{4.12}
\{x \in \mathbb{R} : E_{j}[|f|](x) > \lambda, \text{ 且 } k < j \text{ 时 } E_{k}[|f|](x) \leq \lambda\}
\label{eq:4.12}
\end{equation}
可以分割为 $I_{m}^{j} \in \mathcal{D}_{j}, m = 1, 2, \ldots, j_{M}$ 的和. 让 $j, m$ 变动, 将 $I_{m}^{j}$ 重新排列, 得到 $E_{\lambda}$ 的分割
\begin{equation}\tag{4.13}
E_{\lambda} = \coprod_{j \in J}I_{j}.
\label{eq:4.13}
\end{equation}
设 $x \in I_{j} =: I_{m_{0}}^{j_{0}}$. 由 $I_{m_{0}}^{j_{0}}$ 的取法, $m_{I_{j}}(|f|) = E_{j_{0}}[|f|](x) > \lambda$ 成立. 用 $I_{j}^{*} \in \mathcal{D}_{j_{0}-1}$ 表示二进区间 $I_{j}$ 的父区间, 即包含 $I_{j}$ 的体积为 2 倍的唯一区间. 由 $I_{m_{0}}^{j_{0}}$ 的取法, 即由式~\eqref{eq:4.12}, 必须有 $m_{I_{j}^{*}}(|f|) = E_{j_{0}-1}[|f|](x) \leq \lambda$. 这里,
\begin{align*}
\frac{1}{2}m_{I_{j}}(|f|) &= \frac{1}{2|I_{j}|}\int_{I_{j}}|f(y)|\dd y \\
&= \frac{1}{|I_{j}^{*}|}\int_{I_{j}}|f(y)|\dd y \\
&= \frac{1}{|I_{j}^{*}|}\int_{I_{j}^{*}}|f(y)|\dd y \\
&= m_{I_{j}^{*}}(|f|),
\end{align*}
所以
\[
\frac{1}{2}m_{I_{j}}(|f|) \leq m_{I_{j}^{*}}(|f|) = E_{j_{0}-1}[|f|](x) \leq \lambda,
\]
证明结束.
\end{proof}

由 $\{I_{j}\}_{j \in J} \subset \mathcal{D}$ 互不相交与式~\eqref{eq:4.11} 可知
\begin{equation}\tag{4.14}
|\{x \in \mathbb{R} : M_{\dyadic}f(x) > \lambda\}| = \sum_{j \in J}|I_{j}| \leq \frac{1}{\lambda}\|f\|_{L^{1}}
\label{eq:4.14}
\end{equation}
成立.

下面的定理是奇异积分算子的有界性的关键. $g$ 是 good 的略写, $b$ 是 bad 的略写. 这里的讨论是把第 1.6 节的讨论集大成并应用于函数.

\begin{theorem}[Calderón--Zygmund 分解]\label{thm:4.3}
对 $L^{1}$ 函数 $f$ 与 $\lambda > 0$, 存在满足下列性质的至多可数个二进区间的集族 $\{I_{j}\}_{j \in J}$ 与至多可数个 $L^{1}$ 函数的集族 $\{g\} \cup \{b_{j}\}_{j \in J}$.

(1) $f$ 分解为下面的形式.
\begin{equation}\tag{4.15}
f = g + \sum_{j \in J}b_{j} = g + b.
\label{eq:4.15}
\end{equation}

(2)（$L^{1}$ 条件）$g$ 的 $L^{1}$ 范数小于 $f$ 的.
\begin{equation}\tag{4.16}
\int_{\mathbb{R}}|g(x)|\dd x \leq \int_{\mathbb{R}}|f(x)|\dd x.
\label{eq:4.16}
\end{equation}

(3)（$L^{\infty}$ 条件）$g$ 属于 $L^{\infty}(\mathbb{R})$.
\begin{equation}\tag{4.17}
\text{对几乎所有的 } x \in \mathbb{R}, \quad |g(x)| \leq 2\lambda.
\label{eq:4.17}
\end{equation}

(4)（支集条件）对每个 $j \in J$, $b_{j}$ 的支集包含于 $I_{j}$ 的闭包.
\begin{equation}\tag{4.18}
\operatorname{supp}(b_{j}) \subset \overline{I_{j}}.
\label{eq:4.18}
\end{equation}

(5)（矩条件）对每个 $j$, $b_{j}$ 是可积函数, 且满足
\begin{equation}\tag{4.19}
\int_{\mathbb{R}}b_{j}(x)\dd x = 0.
\label{eq:4.19}
\end{equation}

(6)（水平集 $\{M_{\dyadic}f > \lambda\}$ 的分割）$\{I_{j}\}_{j \in J}$ 互不相交, 且
\begin{equation}\tag{4.20}
\{M_{\dyadic}f > \lambda\} = \coprod_{j \in J}I_{j}
\label{eq:4.20}
\end{equation}
成立.

这个分解称为（高度为 $\lambda$ 的）Calderón--Zygmund 分解.
\end{theorem}

\begin{proof}
取互不相交的二进区间的可数序列 $\{I_{j}\}_{j \in J}$, 满足
\begin{equation}\tag{4.21}
\lambda \leq m_{I_{j}}(|f|) \leq 2\lambda, \quad \{x \in \mathbb{R} : M_{\dyadic}f(x) > \lambda\} = \bigcup_{j \in J}I_{j}.
\label{eq:4.21}
\end{equation}
这个区间的序列的存在性已在定理~\ref{thm:4.2} 中证明. 满足式~\eqref{eq:4.18} 与式~\eqref{eq:4.19} 的 $b_{j}$ 如下定义（练习题~\ref{exer:4.8} 参照）.
\begin{equation}\tag{4.22}
b_{j} \equiv \left(f(x) - m_{I_{j}}(f)\right)\chi_{I_{j}}.
\label{eq:4.22}
\end{equation}
与 $b_{j}$ 相应, 令 $g \equiv f - \sum_{j \in J}b_{j}$, 则至少可以得到 $f$ 的分解. 这里注意, 和至少在"几乎所有的"的意义与 $L^{1}$ 的意义两种意义下成立. 这是因为 $b_{j}$ 的支集不共有内部. 由 $b_{j}$ 的定义, $g$ 的 $L^{1}$ 条件式~\eqref{eq:4.16} 成立. 剩下的只是对 $g$ 证明式~\eqref{eq:4.17}.

设 $x \in I_{j}$. 由于 $\{I_{\tilde{j}}\}_{\tilde{j} \in J}$ 互不相交,
\[
g(x) = f(x) - \sum_{j \in J}b_{j}(x) = f(x) - f(x) + m_{I_{j}}(f) = m_{I_{j}}(f)
\]
成立. 因此, 由式~\eqref{eq:4.21}, $|g(x)| \leq 2\lambda$. 反之, 设 $x \notin \bigcup_{j \in J}I_{j}$. 由 Lebesgue 微分定理, 式~\eqref{eq:4.17} 由 $|g(x)| = |f(x)| \leq M_{\dyadic}f(x) \leq \lambda$ 得出.
\end{proof}

\begin{note}\label{note:4.1}
$f \in L^{1}(\mathbb{R}) \cap L^{2}(\mathbb{R})$ 时, 在 $L^{1}(\mathbb{R}) \cap L^{2}(\mathbb{R})$ 的拓扑下 $b = \sum_{j \in J}b_{j}$, 且 $g$ 属于 $L^{1}(\mathbb{R}) \cap L^{2}(\mathbb{R})$, 这由 $g$ 的构造可知.
\end{note}

\begin{note}\label{note:4.2}
引理~\ref{lem:4.2}, 定理~\ref{thm:4.2}, \ref{thm:4.3} 对 $f \in L^{2}(\mathbb{R})$ 也同样成立.
\end{note}

证明 Hilbert 变换的弱 (1,1) 有界性.

\begin{theorem}\label{thm:4.4}
对一切 $\lambda$ 与 $f \in L^{1}(\mathbb{R}) \cap L^{2}(\mathbb{R})$,
\begin{equation}\tag{4.23}
|\{|Hf| > \lambda\}| \lesssim \frac{1}{\lambda}\|f\|_{L^{1}}
\label{eq:4.23}
\end{equation}
成立.
\end{theorem}

首先, 取 $f \in L^{1}(\mathbb{R}) \cap L^{2}(\mathbb{R})$ 这样的形式, 不只要求 $f \in L^{1}(\mathbb{R})$ 还要求 $f \in L^{2}(\mathbb{R})$, 是因为要使 $Hf$ 这个记号有意义, 眼下需要如此.

\begin{proof}
将 $f$ 按高度 $\lambda$ 作 Calderón--Zygmund 分解, 表为 $f = g + \sum_{j \in J}b_{j} = g + b$. 由注意~\ref{note:4.1} 可知 $b, g \in L^{2}(\mathbb{R})$, 所以 $Hf = Hg + Hb$ 成立. 转到补集考虑显然地,
\begin{equation}\tag{4.24}
\{|Hf| > \lambda\} \subset \left\{|Hg| > \frac{\lambda}{2}\right\} \cup \left\{|Hb| > \frac{\lambda}{2}\right\}
\label{eq:4.24}
\end{equation}
成立. 因此, 只需证明
\begin{equation}\tag{4.25}
\left|\left\{|Hg| > \frac{\lambda}{2}\right\}\right| \lesssim \frac{1}{\lambda}\|f\|_{L^{1}},
\label{eq:4.25}
\end{equation}
\begin{equation}\tag{4.26}
\left|\left\{|Hb| > \frac{\lambda}{2}\right\}\right| \lesssim \frac{1}{\lambda}\|f\|_{L^{1}}
\label{eq:4.26}
\end{equation}
这两个. 对应好函数的式~\eqref{eq:4.25} 容易证明. 由 Chebyshev 不等式,
\[
\left|\left\{|Hg| > \lambda\right\}\right| \leq \frac{1}{\lambda^{2}}\int_{\mathbb{R}}|Hg(x)|^{2}\dd x
\]
成立. 因此, 依次使用 $H$ 的 $L^{2}$ 有界性, 式~\eqref{eq:4.17}, 式~\eqref{eq:4.16},
\[
|\{|Hg| > \lambda\}| \lesssim \frac{1}{\lambda^{2}}\int_{\mathbb{R}}|g(x)|^{2}\dd x \lesssim \frac{1}{\lambda}\|g\|_{L^{1}} \lesssim \frac{1}{\lambda}\|f\|_{L^{1}}
\]
成立. 式~\eqref{eq:4.26} 稍难. 如下证明. 将 $I_{j}$ 表为 $I_{j} = [a_{j}-r_{j}, a_{j}+r_{j})$, 令 $2I_{j} = [a_{j}-2r_{j}, a_{j}+2r_{j})$. 另外, 对这种 $I_{j}$, 令 $c(I_{j}) = a_{j}$. 首先, 由式~\eqref{eq:4.14} 与式~\eqref{eq:4.20},
\begin{equation}\tag{4.27}
\left|\bigcup_{j \in J}2I_{j}\right| \leq 2\sum_{j \in J}|I_{j}| = 2|\{M_{\dyadic}f > \lambda\}| \leq \frac{2}{\lambda}\|f\|_{L^{1}},
\label{eq:4.27}
\end{equation}
所以只需证明
\begin{equation}\tag{4.28}
\left|\{|Hb| > \lambda\} \setminus \bigcup_{j \in J}2I_{j}\right| \lesssim \frac{1}{\lambda}\|f\|_{L^{1}},
\label{eq:4.28}
\end{equation}
但式~\eqref{eq:4.28} 的左边由 Chebyshev 不等式（定理~\ref{thm:2.17}）$\leq \frac{1}{\lambda}\sum_{j \in J}\int_{\mathbb{R}\setminus 2I_{j}}|Hb_{j}(x)|\dd x$, 所以式~\eqref{eq:4.26} 归结为
\begin{equation}\tag{4.29}
\int_{\mathbb{R}\setminus 2I_{j}}|Hb_{j}(x)|\dd x \lesssim \int_{I_{j}}|b_{j}(x)|\dd x \quad (j \in J).
\label{eq:4.29}
\end{equation}
设 $x \in \mathbb{R}\setminus 2I_{j}$, $y \in I_{j}$, 注意
\[
2|x-y| \geq \ell(I_{j}) \geq |y-c_{j}|
\]
且
\[
\left|\frac{1}{x-y} - \frac{1}{x-c_{j}}\right| \leq \frac{|y-c_{j}|}{|(x-y)(x-c_{j})|} \leq \frac{2\ell(I_{j})}{|x-c_{j}|^{2}}.
\]
$x \in \mathbb{R}\setminus 2I_{j}$ 时, 由这个逐点估计与关于 $b_{j}$ 的支集条件式~\eqref{eq:4.18} 和矩条件式~\eqref{eq:4.19}, 制造正负相消,
\[
Hb_{j}(x) = \int_{I_{j}}\frac{1}{x-y}b_{j}(y)\dd y = \int_{I_{j}}\left(\frac{1}{x-y} - \frac{1}{x-c(I_{j})}\right)b_{j}(y)\dd y,
\]
对其使用三角不等式,
\[
|Hb_{j}(x)| \leq \int_{I_{j}}\left|\frac{1}{x-y} - \frac{1}{x-c(I_{j})}\right| \cdot |b_{j}(y)|\dd y \lesssim \frac{\ell(I_{j})^{2}m_{I_{j}}(|b_{j}|)}{|x-c(I_{j})|^{2}},
\]
所以将其代入式~\eqref{eq:4.29}, 在 $\mathbb{R}\setminus 2I_{j}$ 上积分即可. 即, 将正负相消巧妙地累加即可.
\end{proof}

\begin{exercise}\label{exer:4.6}
将 $\mathcal{D}_{0}(\mathbb{R})$, $\mathcal{D}_{1}(\mathbb{R})$, $\mathcal{D}_{2}(\mathbb{R})$ 在数直线 $\mathbb{R}$ 上图示.
\end{exercise}

\begin{exercise}\label{exer:4.7}
对 $\mathbb{R}$ 上的可测函数 $f$,
\[
M_{\dyadic}f = \chi_{(-\infty, 0)}M_{\dyadic}[\chi_{(-\infty, 0)}f] + \chi_{[0, \infty)}M_{\dyadic}[\chi_{[0, \infty)}f]
\]
成立, 证明之.
\end{exercise}

\begin{exercise}\label{exer:4.8}
设 $f$ 是局部可积函数, $Q$ 是立方体. 令
\[
B(x) = \chi_{Q}(x)f(x) - \frac{\chi_{Q}(x)}{|Q|}\int_{Q}f(z)\dd z,
\]
证明 $B$ 是可积函数, 且满足 $\int_{Q}B(x)\dd x = 0$.
\end{exercise}

\begin{exercise}\label{exer:4.9}
设 $E$ 是具有有限测度的 Lebesgue 可测集合, 且能表为 $\mathcal{D}(\mathbb{R})$ 的元的并. 即, 存在 $\mathcal{A} \subset \mathcal{D}(\mathbb{R})$, 使得
\[
E = \bigcup_{Q \in \mathcal{A}}Q.
\]
(1) 证明 $\sup_{Q \in \mathcal{A}}|Q| < \infty$ 成立.

(2) 证明存在 $\mathcal{A}_{0} \subset \mathcal{A}$, 使得 $\sum_{Q \in \mathcal{A}_{0}}\chi_{Q} = \chi_{E}$ 成立.
\end{exercise}

\begin{exercise}\label{exer:4.10}
计算 $t > 2$ 时的 $H\chi_{[-1,1]}(t)$, 证明 $H\chi_{[-1,1]} \notin L^{1}([2, \infty))$.
\end{exercise}
